\documentclass[11pt]{article}
\usepackage[a4paper,margin=1in]{geometry}
\usepackage{amsmath,amssymb,booktabs,xcolor,enumitem}
\usepackage[colorlinks=true,linkcolor=blue,urlcolor=blue]{hyperref}
\usepackage{listings}
\lstset{basicstyle=\ttfamily\small,columns=fullflexible,breaklines=true,frame=single,
        xleftmargin=1em,framexleftmargin=1em}
\setlength{\parindent}{0pt}\setlength{\parskip}{5pt}
\newcommand{\state}[1]{\textsc{#1}}
\newcommand{\code}[1]{\texttt{#1}}
\title{\textbf{Device-Simulation GUI --- Architecture Spec v0.2.9 Delta}\\[3pt]
\large Plateau-Measurement-Integrity Contract, \state{needs\_recheck} State, and the
Audit-Correction / Supersession Pattern}
\date{2026-08-23}\author{}
\begin{document}\maketitle\vspace{-1.5em}

\noindent\fcolorbox{black}{yellow!12}{\parbox{0.97\linewidth}{\small
\textbf{Scope.} This delta specifies three items flagged for v0.2.9 that are not yet in the v0.2.8/v0.2.9
master: (1) the \emph{Plateau-Measurement-Integrity Contract} (PMIC); (2) the explicit
\state{needs\_recheck} verification state and its transitions; (3) the audit-correction / supersession
pattern for defective-but-preserved evidence. It is written to be folded into the master spec's
verification-checkpoint chapter. It is grounded in a concrete defect resolved in the backend
(Liang-1997 SHE-BTE reproduction, 71$\times$58 node, 2026-08-23); that case is used as the worked
example so the contract is testable, not aspirational.}}

\section{Motivating defect (normative example)}
A convergence observable (rate $G$) was accepted as ``plateaued'' when the relative spread of its last
four values fell below a gate. The solver stopped each continuation \emph{chunk} on its own residual
criterion \code{tol\_phi}. When the solution was already near its fixed point, a chunk stopped after a
\emph{single} iteration; the ``last four'' then spanned one value, whose spread is identically $0$ ---
a \textbf{degenerate false pass}. The physical value was fine; the \emph{measurement of acceptance} was
not. The GUI must make this class of error unrepresentable.

\section{Plateau-Measurement-Integrity Contract (PMIC)}
\textbf{Intent.} A plateau verdict may be computed only over genuine, independent samples of the
observable, accumulated across solver invocations, and never manufactured by an early solver stop.

\subsection{Definitions}
\begin{description}[leftmargin=1.4em,style=nextline]
\item[Real iteration] one outer fixed-point update that actually advanced the observable; identified by
a strictly increasing absolute iteration index \code{abs\_iter} on the node's cumulative trajectory.
\item[Cumulative window] the last $W{=}4$ \emph{real} iterations of the node's trajectory, taken across
chunk boundaries --- \emph{not} the last four of any single chunk.
\item[Rolling buffer] the persisted, ordered list of the most recent $\le W$ real samples
$\{(\code{abs\_iter},\,\text{observables})\}$, carried between solver invocations in the node sidecar.
\item[\code{plateau\_evaluable}] a boolean: \textbf{false} until the buffer holds $W$ genuine samples.
\end{description}

\subsection{Normative rules (MUST)}
\begin{enumerate}[leftmargin=1.6em]
\item The plateau gate is evaluated over the cumulative window, not a chunk-local window.
\item \code{plateau\_evaluable} is \textbf{false} while $\lvert\text{buffer}\rvert < W$; a node MUST NOT
be certified while not evaluable, \emph{even if} a chunk-local spread would be $0$.
\item The solver-stopping residual (\code{tol\_phi}) is \textbf{decoupled} from acceptance: it may end a
chunk, but MUST NOT short-circuit or stand in for the observable plateau decision.
\item The production spread statistic and thresholds are \textbf{fixed} and identical to the backend:
$\operatorname{spread}(x)=(\max x-\min x)/\bar x$, gate $\le 3\times10^{-4}$ on \emph{both} rate
observables, residual floored $\le 1.5\times10^{-3}$. The GUI MUST source these from one shared
definition, never redefine them.
\item \textbf{Buffer integrity (fail-closed).} On resume, \code{abs\_iter} across the buffer MUST be
strictly increasing (hence duplicate-free). A violation aborts certification (a duplicated index could
fabricate an artificial plateau); it is a hard error, not a warning.
\item \textbf{Provenance.} Each certification records \code{buffer\_origin} $\in$
\{\code{persisted\_sidecar}, \code{npz\_terminal\_reconstruction}, \code{empty\_fresh}\}, so the evidence
path that certified the gate is auditable after the fact.
\item A spread of $0$ is legitimate \emph{iff} \code{plateau\_evaluable} and the window holds $W$ distinct
real iterations (e.g.\ four independent re-solves re-confirming one fixed point); it is forbidden
otherwise.
\end{enumerate}

\subsection{Contract surface (data the GUI persists per node)}
\begin{lstlisting}
sidecar.plateau = {
  buffer:           [ {abs_iter, <observables...>, residual}, ... up to W ],
  plateau_evaluable: bool,          # false until |buffer| == W
  passed:            bool,          # gate over the cumulative window
  buffer_origin:     enum,          # persisted_sidecar | npz_terminal_reconstruction | empty_fresh
  window_iters:      [abs_iter...], # the exact indices used (independently checkable)
}
\end{lstlisting}

\subsection{Checkpoint hook}
The plateau checkpoint (the master's convergence-ladder rung for a node) reports \state{pass} only when
\code{plateau\_evaluable $\wedge$ passed}; it reports \state{not\_evaluable} (a distinct, non-failing,
non-passing status) while the window is short; and it reports \state{fail} when evaluable and the gate is
not met. \state{not\_evaluable} MUST drive ``continue accumulating'', never ``accept''.

\section{The \state{needs\_recheck} verification state}
\textbf{Intent.} A node whose \emph{acceptance evidence} is invalid --- as opposed to a node that has
genuinely failed or genuinely passed --- occupies a first-class state that blocks downstream use without
asserting a scientific failure.

\subsection{State set and transitions}
\begin{center}\small
\begin{tabular}{lll}\toprule
from & event & to \\ \midrule
\state{accumulating} & \code{plateau\_evaluable $\wedge$ passed} & \state{certified} \\
\state{accumulating} & evaluable $\wedge\,\neg$passed, cap reached & \state{failed} \\
\state{certified} & audit finds evidence degenerate/ambiguous & \state{needs\_recheck} \\
\state{needs\_recheck} & recheck yields evaluable genuine window, passed & \state{certified} \\
\state{needs\_recheck} & recheck evaluable $\wedge\,\neg$passed & \state{failed} \\
\state{needs\_recheck} & cannot establish a valid window under policy & \state{audit\_incomplete} \\
\bottomrule\end{tabular}
\end{center}

\subsection{Normative rules (MUST)}
\begin{enumerate}[leftmargin=1.6em]
\item \state{needs\_recheck} is neither \state{pass} nor \state{fail}; the GUI MUST render it as its own
badge and MUST block any campaign-level verdict (e.g.\ a GREEN/YELLOW classification, or an
extrapolation) that depends on the node until it leaves this state.
\item Entering \state{needs\_recheck} MUST NOT overwrite the original recorded metrics row; the original
(defective) record is preserved verbatim (see \S\ref{sec:supersede}).
\item A recheck MUST reuse the identical run definition (mesh, grids, physics, solver settings,
thresholds) and change \emph{only} the acceptance measurement; it carries a new \code{run\_id} and names
the prior checkpoint as \code{parent}. A pure safety cap (max continuation chunks) may be raised and MUST
be labeled as such, never as a scientific-gate change.
\item Reconstruction subtlety: a short terminal chunk does \emph{not} automatically imply
\state{needs\_recheck}. If $\ge W$ genuine cumulative samples are retrievable (persisted buffer, or a
full-window terminal reconstruction) the node is evaluable; \state{needs\_recheck} is entered only when a
valid window is unretrievable or ambiguous.
\end{enumerate}

\section{Audit-correction / supersession pattern}\label{sec:supersede}
\textbf{Intent.} Corrections are additive and immutable: nothing that was recorded is edited or deleted;
a superseding record points back to what it replaces, and the chain is independently verifiable by hash.

\subsection{Normative rules (MUST)}
\begin{enumerate}[leftmargin=1.6em]
\item \textbf{Preserve.} The original artifact (data checkpoint, metrics row, generated report) is kept
byte-for-byte under an archival name; its content hash is recorded.
\item \textbf{Supersede, don't overwrite.} A new \emph{audit-correction record} is emitted as a separate
artifact carrying: \code{original\_status}, \code{audit\_status}, \code{reason}, the superseded
artifact's \code{sha256}, the new \code{run\_id}/\code{parent}, and the corrected decision. The record
itself gets a \code{sha256}.
\item \textbf{Derive, don't mutate.} A reconstructed or corrected decision is emitted as its own artifact;
source data files are never rewritten to ``fix'' a metric.
\item \textbf{Label re-derived verdicts.} A classification recomputed after a measurement repair MUST be
visibly labeled as re-derived (e.g.\ ``YELLOW --- re-derived after plateau-bookkeeping repair'') so it
cannot be confused with the earlier defective output. Reports MUST present \emph{scientific reproduction}
and \emph{numerical parity} as separate statements.
\item \textbf{Lineage.} If a corrected node was a warm-start parent of another accepted node, the record
states whether the corrected value moved materially; if it did not, dependent nodes are \emph{not}
regenerated, and the record says so explicitly.
\end{enumerate}

\subsection{Audit-correction record (schema)}
\begin{lstlisting}
audit_correction = {
  node, original_status, audit_status,   # e.g. PASS -> NEEDS_RECHECK
  reason,                                # e.g. "degenerate_plateau_window"
  superseded_artifact, superseded_sha256,
  run_id, parent_checkpoint_sha256,
  resolution: { new_status, evidence, window_iters, buffer_origin },
  sha256                                 # of this record
}
\end{lstlisting}

\section{Worked example (normative acceptance test for the GUI)}
The GUI's verification layer MUST reproduce this sequence for the 71$\times$58 node:
\begin{enumerate}[leftmargin=1.6em,itemsep=1pt]
\item Original acceptance: terminal chunk ran $1$ real iteration; chunk-local spread $=0$ $\Rightarrow$
\emph{defective} \state{certified}. Audit reclassifies to \state{needs\_recheck}; original row preserved;
audit-correction record emitted (\code{reason=degenerate\_plateau\_window}).
\item Recheck (new \code{run\_id}, same run definition, safety cap raised only): the node early-stops
$1$ real iter/chunk. PMIC holds \state{not\_evaluable} at cumulative $n{=}1,2,3$ (indices $359,360,361$)
and certifies at $n{=}4$ (index $362$): four distinct re-solves, both spreads $=0$ ($G_{ii}$ spread $=0$;
$G_{\mathrm{tot}}$ spread $=0$), residual floored.
\state{needs\_recheck} $\rightarrow$ \state{certified}.
\item Campaign verdict recomputed on the certified nodes; classification labeled ``re-derived after
plateau-bookkeeping repair''; scientific reproduction (significant non-monotone reversal) and numerical
parity ($\pm2\%$ archival) reported separately; lineage note records the parent value was unchanged, so
the dependent node is not regenerated.
\end{enumerate}
Any GUI build that would let step~1 stand as a pass, that mutates the original row in step~1, or that
omits the re-derived label / parity separation in step~3, \textbf{fails} this acceptance test.

\end{document}
